% Abstract (English and Russian) and the statement on AI tools. Finalise after the pilot:
% the results sentence is marked \pending until RQ-1…RQ-3 have numbers with a passport (docs/thesis/figures.md).
\chapter*{Abstract}
\addcontentsline{toc}{chapter}{Abstract}
\thispagestyle{plain}

Adults who learn to read the Qur'an usually start with a primer and a teacher who listens to every
syllable. This thesis asks how much of that listening can be done by a machine for Russian-speaking
beginners, and builds the instrument needed to answer the question. The instrument is a Telegram tutor
that follows the classical primer \emph{Muallim al-Sani} by Ahmad Hadi Maksudi in the book's own order:
125 lessons and 800 cards, each traceable to a page and line of the 2010 edition. The learner sends a
voice message; the open model Quran Muaalem recognises phonemes; a rule-based alignment classifies
differences into ten error kinds tied to syllables; and the learner receives a short explanation in
Russian with a reference recording. Every attempt keeps a full decision trace, so that verdicts can be
recomputed later.

The thesis follows design science research and evaluates two artefacts. The tutor is evaluated by a
measurement study of agreement with two teachers (RQ-1), an analysis of beginners' errors that separates
learner errors from model errors (RQ-2), and a within-subject pilot experiment comparing full
explanations with a bare correct/incorrect signal (RQ-3). The development harness --- hooks that freeze
tests, protect human-owned zones and stop the agent while checks fail, together with reviewer sub-agents
and journals --- is evaluated as an embedded case study (RQ-4). Hypotheses and the analysis plan are
preregistered.

Preliminary checks show that the pipeline separates correct from incorrect expert recitation (99--100
versus 40--67 points) and answers within the 10-second target. \pending{pilot results for RQ-1--RQ-3}.
The contribution is the first teacher-referenced evaluation of an open Qur'anic pronunciation model on
primer-level reading by second-language beginners, a consented, doubly annotated dataset, and an
evidence-based account of building research software with AI coding agents under mechanical
guard-rails.

\medskip
\noindent\textbf{Keywords:} computer-assisted pronunciation training; mispronunciation detection;
Qur'anic Arabic; Telegram bot; design science; AI coding agents; harness engineering.

\chapter*{\textrussian{Аннотация}}
\addcontentsline{toc}{chapter}{Аннотация (Russian abstract)}
\thispagestyle{plain}
\begin{otherlanguage}{russian}
Взрослые, которые учатся читать Коран, обычно начинают с букваря и учителя, слушающего каждый слог.
В диссертации исследуется, какую часть этого слушания может взять на себя машина для русскоязычных
начинающих, и создаётся инструмент для ответа на этот вопрос --- Telegram-тренажёр, который ведёт
ученика по букварю Ахмад-Хади Максуди «Муаллим ас-Сани» в порядке самой книги: 125 уроков и 800
карточек, каждая прослеживается до страницы и строки издания 2010~г. Ученик присылает голосовое;
открытая модель Quran Muaalem распознаёт фонемы; выравнивание по правилам относит расхождения к
десяти типам ошибок с привязкой к слогам; ученик получает короткое объяснение на русском и эталонную
запись. По каждой попытке сохраняется полный след решения.

Работа выполнена в парадигме design science и оценивает два артефакта. Тренажёр оценивается
измерением согласия с двумя преподавателями (RQ-1), анализом ошибок начинающих с отделением ошибок
ученика от ошибок модели (RQ-2) и внутрисубъектным пилотным экспериментом, сравнивающим полные
объяснения с сигналом «верно / неверно» (RQ-3). Обвязка разработки --- хуки, замораживающие тесты,
защищающие зоны человека и не дающие агенту закончить при красных проверках, вместе с
субагентами-ревьюерами и журналами --- оценивается как встроенное кейс-стади (RQ-4). Гипотезы и план
анализа предрегистрированы.

Предварительные проверки показывают, что конвейер отличает верное чтение эксперта от неверного
(99--100 против 40--67 баллов) и отвечает быстрее 10 секунд. \pending{результаты пилота по RQ-1--RQ-3}.
Вклад работы --- первая оценка открытой модели проверки чтения Корана на букварном чтении начинающих,
для которых арабский не родной, относительно оценки преподавателей; набор данных с двойной разметкой,
собранный с согласия участников; доказательное описание разработки исследовательского ПО ИИ-агентами
под механическим контролем.

\medskip
\noindent\textbf{Ключевые слова:} компьютерное обучение произношению; обнаружение ошибок
произношения; чтение Корана; Telegram-бот; design science; ИИ-агенты; обвязка агента.
\end{otherlanguage}

\chapter*{Statement on the use of AI tools}
\addcontentsline{toc}{chapter}{Statement on the use of AI tools}
\thispagestyle{plain}

The software described in this thesis was developed by the author together with an AI coding agent
(Claude Code). The agent wrote most of the code and drafts of documentation under the rules described in
Section~\ref{sec:impl-harness}; every test was accepted by the author before implementation, and all
decisions on teaching content, texts shown to learners, thresholds, consent and data were taken by the
author. Drafts of this text were prepared with the same agent from the project's own records and revised
by the author. Every prompt to the agent is stored with its date in the project journal, which makes the
use of AI auditable. The agent did not annotate data and did not compute any reported result outside the
versioned code that is cited for each table. Bibliographic data were checked against Crossref, arXiv or
the publisher. This statement will be aligned with the university's policy on AI use once it is
confirmed.
